\documentclass[11pt,a4paper]{article}
\usepackage[T1]{fontenc}
\usepackage[utf8]{inputenc}
\usepackage[french]{babel}
\usepackage{amsmath,amssymb,geometry,multicol,array,enumitem,xcolor}
\geometry{margin=1.8cm}
\setlength{\parindent}{0pt}
\setlength{\parskip}{5pt}
\newcommand{\rep}[1][1.1cm]{\par\vspace{#1}}
\newcommand{\dd}{\mathrm{dB}}
\newcommand{\jw}{\mathrm{j}\omega}
\newcommand{\ex}[2]{\subsection*{Exercice #1 --- #2}}
\begin{document}
\begin{center}
{\Large\bfseries Fiche d'entraînement aux calculs}\par
{\large Étude harmonique et diagrammes de Bode}\par
%\small Nom : \rule{6cm}{0.4pt}\hfill Date : \rule{3cm}{0.4pt}
\end{center}
\vspace{-2mm}
\fbox{\begin{minipage}{0.96\linewidth}
\textbf{Outils à maîtriser} :
\(\mathrm{j}^2=-1\), \( |a+\mathrm{j}b|=\sqrt{a^2+b^2}\),
\(\arg(a+\mathrm{j}b)=\operatorname{arctan}\left(\dfrac{b}{a}\right)\),
\(G_{\dd}=20\log_{10}|H(\jw)|\),
\(\omega_c=1/\tau\).
Pour un produit, les gains en dB et les phases s'additionnent ; pour un quotient, ils se soustraient.
\end{minipage}}

\ex{1}{Puissances de dix et logarithmes (5 min)}
Sans calculatrice, compléter :
\begin{center}\begin{tabular}{|c|c|c|c|c|}\hline
$x$ & $0{,}01$ & $0{,}1$ & $10$ & $100$ \\ \hline
$\log_{10}(x)$ & \rule{0pt}{0.8cm}\hspace{1.4cm} & \hspace{1.4cm} & \hspace{1.4cm} & \hspace{1.4cm} \\ \hline
$20\log_{10}(x)$ (dB) & \rule{0pt}{0.8cm} & & & \\ \hline
\end{tabular}\end{center}
Calculer ensuite à la calculatrice : \(20\log_{10}(2)\), \(20\log_{10}(1/\sqrt{2})\), \(20\log_{10}(5)\) (arrondir à 0,1 dB).

\ex{2}{Nombres complexes (8 min)}
Pour chacun des nombres suivants, calculer le module et l'argument en degrés :
\begin{multicols}{2}
\begin{enumerate}[label=\alph*)]
\item $z_1=3+4\mathrm{j}$
\item $z_2=1+\mathrm{j}$
\item $z_3=-1+\mathrm{j}$
\item $z_4= -2\mathrm{j}$
\end{enumerate}
\end{multicols}
\textit{Attention : le signe de la partie réelle détermine le quadrant.}

\ex{3}{Passer d'une fonction de transfert à $H(\mathrm{j}\omega)$ (5 min)}
Remplacer $p$ par $\mathrm{j}\omega$ puis simplifier en utilisant $\mathrm{j}^2=-1$ :
\begin{enumerate}[label=\alph*)]
\item $H_1(p)=\dfrac{4}{1+0{,}2p}$
\item $H_2(p)=\dfrac{3}{p}$
\item $H_3(p)=\dfrac{2}{1+0{,}4p+0{,}04p^2}$
\end{enumerate}

\ex{4}{Gain pur et intégrateur (8 min)}
\begin{enumerate}[label=\alph*)]
\item Pour $H(p)=5$, déterminer $|H(\jw)|$, le gain en dB et la phase.
\item Pour $H(p)=\dfrac{10}{p}$, compléter le tableau (\(\omega\) en rad/s) :
\end{enumerate}
\begin{center}\begin{tabular}{|c|c|c|c|}\hline
$\omega$ & $0{,}1$ & $1$ & $10$ \\ \hline
$|H(\jw)|$ & \rule{0pt}{0.85cm} & & \\ \hline
$G_{\dd}$ & \rule{0pt}{0.85cm} & & \\ \hline
$\varphi$ & \rule{0pt}{0.85cm} & & \\ \hline
\end{tabular}\end{center}
Quelle est la variation du gain lorsqu'on multiplie la pulsation par 10 ? \rep[0.8cm]

\ex{5}{Système du premier ordre (12 min)}
On étudie \(H(p)=\dfrac{4}{1+0{,}1p}\).
\begin{enumerate}[label=\alph*)]
\item Identifier $K$ et $\tau$, puis calculer la pulsation de cassure $\omega_c$.
\item Montrer que \( |H(\jw)|=\dfrac{4}{\sqrt{1+(0{,}1\omega)^2}}\) et \(\varphi=-\arctan(0{,}1\omega)\).
\item Compléter le tableau (arrondir les gains à 0,1 dB et les phases au degré).
\end{enumerate}
\begin{center}\begin{tabular}{|c|c|c|c|}\hline
$\omega$ (rad/s) & $1$ & $10$ & $100$ \\ \hline
$|H(\jw)|$ & \rule{0pt}{0.9cm} & & \\ \hline
$G_{\dd}$ & \rule{0pt}{0.9cm} & & \\ \hline
$\varphi$ (degrés) & \rule{0pt}{0.9cm} & & \\ \hline
\end{tabular}\end{center}
\begin{enumerate}[label=\alph*),start=4]
\item À $\omega=\omega_c$, comparer le gain en dB à sa valeur à basse fréquence.
\end{enumerate}

\ex{6}{Second ordre : calcul direct (12 min)}
On considère \(H(p)=\dfrac{2}{1+0{,}4p+0{,}04p^2}\).
\begin{enumerate}[label=\alph*)]
\item Mettre le dénominateur de $H(\jw)$ sous la forme $a(\omega)+\mathrm{j}b(\omega)$.
\item Pour $\omega=5$ rad/s, calculer $a$, $b$, $|H(\jw)|$, $G_{\dd}$ et $\varphi$.
\item Vérifier que $\omega_0=5$ rad/s et que la phase vaut $-90^\circ$ à cette pulsation.
\end{enumerate}
\textit{Pour l'argument d'un complexe, utiliser la fonction atan2 ou contrôler le quadrant.}

\ex{7}{Défi : second ordre sous-amorti (10 min)}
On donne \(H(p)=\dfrac{1}{1+\frac{0{,}4}{10}p+\frac{p^2}{100}}\), soit $K=1$, $z=0{,}2$ et $\omega_0=10$ rad/s.
\begin{enumerate}[label=\alph*)]
\item À $\omega=\omega_0$, calculer le module, le gain en dB et la phase.
\item Le gain à cette pulsation est-il supérieur ou inférieur au gain statique ? De combien de dB ?
\end{enumerate}

\newpage
\section*{Corrigé synthétique (à retirer avant distribution)}
\textbf{1.} $\log_{10}(x)$ : $-2,-1,1,2$ ; gains : $-40,-20,20,40$ dB. Puis $6{,}0$ dB ; $-3{,}0$ dB ; $14{,}0$ dB.

\textbf{2.} $z_1$ : $5$, $53{,}1^\circ$ ; $z_2$ : $\sqrt2$, $45^\circ$ ; $z_3$ : $\sqrt2$, $135^\circ$ ; $z_4$ : $2$, $-90^\circ$.

\textbf{3.} a) $4/(1+0{,}2\mathrm{j}\omega)$ ; b) $3/(\mathrm{j}\omega)=-3\mathrm{j}/\omega$ ; c) $2/(1-0{,}04\omega^2+0{,}4\mathrm{j}\omega)$.

\textbf{4.} a) Module $5$, gain $14{,}0$ dB, phase $0^\circ$. b) Modules $100,10,1$ ; gains $40,20,0$ dB ; phases $-90^\circ$ partout. Pente $-20$ dB/décade.

\textbf{5.} $K=4$, $\tau=0{,}1$ s, $\omega_c=10$ rad/s. Pour $\omega=1,10,100$ : modules $3{,}980;2{,}828;0{,}398$ ; gains $12{,}0;9{,}0;-8{,}0$ dB ; phases $-6^\circ;-45^\circ;-84^\circ$. À la cassure : $-3{,}0$ dB par rapport au plateau $20\log_{10}4\simeq12{,}0$ dB.

\textbf{6.} $a=1-0{,}04\omega^2$, $b=0{,}4\omega$. À $\omega=5$ : $a=0$, $b=2$, module $1$, gain $0$ dB, phase $-90^\circ$. Forme canonique : $\omega_0=5$ rad/s, $z=1$.

\textbf{7.} À $\omega=10$ : dénominateur $0+0{,}4\mathrm{j}$ ; module $2{,}5$, gain $20\log_{10}(2{,}5)\simeq 8{,}0$ dB, phase $-90^\circ$. Gain supérieur de $8{,}0$ dB au gain statique ($0$ dB).
\end{document}
